\documentclass{article}

% ============================================================
% Packages
% ============================================================

\usepackage[utf8]{inputenc}
\usepackage[T1]{fontenc}

\usepackage{graphicx}
\usepackage{amsmath}
\usepackage{amssymb}
\usepackage{booktabs}
\usepackage{array}
\usepackage{multirow}
\usepackage{adjustbox}
\usepackage{caption}
\usepackage{natbib}
\usepackage[a4paper,margin=1in]{geometry}

% ============================================================
% Title
% ============================================================

\title{Midterm Thesis Document}
\author{Wassim Mezghanni}
\date{May 2026}

\begin{document}

\maketitle

% ============================================================
% SECTION 0: Thesis Title & Goal
% ============================================================

\section{Thesis Title and Goal}
\label{sec:goal}

\textbf{Title:} Cross-Lingual Alignment as a Proxy for Multilingual LLM Performance

\medskip

The goal of this thesis is to use the evaluation framework proposed by
\citet{kargaran2025mexa} and apply it to a broader set of large language
models (LLMs), including the Llama, Mistral, Qwen3, and Apertus families.

MEXA estimates multilingual coverage without expensive downstream evaluations
by measuring how well non-English sentence representations align with English
representations across the hidden layers of decoder-only LLMs.

The thesis validates the replication on two parallel corpora
(FLORES-200 and the Super Parallel Bible Corpus) against three downstream
benchmarks (Belebele, m-MMLU, and m-ARC), and extends the analysis to
low-resource languages and new model architectures not covered by the
original paper.

% ============================================================
% SECTION 1: Methodology
% ============================================================

\section{Methodology}
\label{sec:methodology}

The core methodology replicates and extends the MEXA framework
proposed by \citet{kargaran2025mexa}.

\paragraph{Theoretical foundation.}

Autoregressive English-centric LLMs use English as an internal semantic pivot
in their intermediate layers \citep{wendler2024llamas,zhao2024multilingualism}.
Measuring how well non-English representations align with English at each
layer therefore acts as a proxy for multilingual task performance.

\paragraph{Alignment formulation.}

For parallel sentences in a target language $L_1$ and English $L_2$,
sentence embeddings $e_{L_1}^{(l)}$ and $e_{L_2}^{(l)}$ are extracted at
each transformer layer $l$, forming a cosine similarity matrix $C$ of size
$n \times n$.

To overcome anisotropy and hubness, MEXA uses a bidirectional retrieval
criterion:

\begin{equation}
\mu C(L_1, L_2, m, l)
=
\frac{1}{n}
\sum_{i=1}^{n}
\mathbb{1}
\left(
c_{ii}(l)
>
\max_{j \neq i}
\{ c_{ij}(l),\, c_{ji}(l) \}
\right)
\end{equation}

where $c_{ij}(l)$ is the cosine similarity between target sentence $i$
and English sentence $j$ at layer $l$.

\paragraph{Embedding extraction.}

Two strategies are used:

\begin{itemize}
    \item \textbf{Last Token}: the hidden state of the final input token.
    
    \item \textbf{Position-Weighted Average}:
\end{itemize}

\begin{equation}
e_l
=
\sum_{t=1}^{T}
w_t h_t^{(l)},
\qquad
w_t
=
\frac{t}{\sum_{k=1}^{T} k}
\end{equation}

This weighting mitigates causal attention bias toward early tokens.

\paragraph{Layer-wise aggregation.}

Scores are summarised using:

\begin{itemize}
    \item $\mu_{\text{Max}}$: peak alignment across layers.
    \item $\mu_{\text{Mean}}$: average alignment across layers.
\end{itemize}

\paragraph{Datasets.}

\textbf{FLORES-200}: a 116-language subset
(100 sentences overlapping with Belebele)
and the full 204-language devtest set
(2,000 sentences).

\textbf{Super Parallel Bible Corpus (sPBC)}:
a 101-language subset (103 verses) for replication
and the full 1,401-language dataset for low-resource analysis.

\paragraph{Models.}

\begin{itemize}
    \item Llama 3.1 8B
    \item Mistral 7B v0.3
    \item Qwen3 8B Base
    \item Qwen3.5 9B Base
    \item Apertus 8B
\end{itemize}

% ============================================================
% SECTION 2: Literature Review
% ============================================================

\section{Literature Review}
\label{sec:literature_review}

The following works form the theoretical and methodological foundation
of this thesis.

\paragraph{\citet{kargaran2025mexa} --- MEXA}

MEXA estimates multilingual LLM coverage by measuring whether
non-English sentence embeddings retrieve their English counterparts
at each transformer layer using a bidirectional retrieval criterion.

Validated on nine LLMs and three downstream benchmarks,
the framework achieves a mean Pearson correlation of $0.90$
with actual multilingual task performance.

This paper forms the central methodological foundation of the thesis.

\paragraph{\citet{wendler2024llamas} --- Do Llamas Work in English?}

Using the logit-lens technique, this work shows that Llama models
internally process non-English inputs through English representations
in middle transformer layers before generating target-language tokens.

This pivot-language hypothesis motivates measuring alignment with
English as a multilingual capability proxy.

\paragraph{\citet{hammerl2024survey} --- Cross-Lingual Alignment Survey}

Provides a taxonomy of multilingual alignment methods, including
contrastive learning, parallel-data alignment objectives,
and post-hoc transformations.

The survey situates MEXA within broader multilingual NLP research.

\paragraph{\citet{hammerl2023anisotropy} --- Anisotropy in Multilingual Models}

Shows that anisotropy and outlier dimensions in multilingual embedding
spaces cause unrelated sentences to appear artificially similar.

This motivates the comparative retrieval-based metric used in MEXA.

\paragraph{\citet{heffernan2022laser3} --- LASER3 / Bitext Mining}

Introduces multilingual sentence encoders evaluated on FLORES-200
for cross-lingual retrieval.

The work is relevant because FLORES-200 is also used in this thesis.

\paragraph{\citet{bandarkar2024belebele} --- Belebele Benchmark}

Introduces a multilingual reading comprehension benchmark
covering 122 language variants.

Belebele is  used to validate MEXA scores.

\paragraph{\citet{nllbteam2022} --- FLORES-200}

Introduces FLORES-200, a human-translated multilingual benchmark
covering 204 language-script pairs.

This dataset is one of the primary datasets used in the thesis.

\paragraph{\citet{dubey2024llama3} --- Llama 3}

Describes the Llama 3 and Llama 3.1 model families.


\paragraph{\citet{rajaee2022isotropy} --- Isotropy in Multilingual BERT}

Studies isotropy variations across multilingual embedding spaces
and explains why comparative metrics are preferable to raw cosine similarity.

\paragraph{\citet{zhao2024multilingualism} --- Multilingual Reasoning in LLMs}

Shows that multilingual LLMs often convert multilingual representations
into English-centric latent representations during reasoning.

\paragraph{\citet{mayer2014bible} --- Parallel Bible Corpus}

Introduces a massively parallel corpus covering more than
1,400 language-script pairs.

The derived Super Parallel Bible Corpus is used in this thesis
for low-resource multilingual evaluation.



% ============================================================
% SECTION 3: Experimental Results
% ============================================================

\section{Experimental Results}
\label{sec:results}

Table~\ref{tab:mexa_results} reports MEXA alignment scores
for all experimental configurations.

Our replication of Llama 3.1 8B  and Mistral 0.3 7B closely matches the paper's
reference values, confirming successful replication.

Qwen3.5 9B achieves the highest scores across most configurations,
suggesting stronger multilingual representations despite similar
parameter counts. In particular, a fine-grained language-level analysis for Qwen3.5 9B Base shows that English (\texttt{eng\_Latn}) consistently achieves the highest alignment score ($\mu_{\text{Max}} = 1.0$) across both the FLORES-200 (116 lang · 100 sent) and Bible (101 lang · sPBC) subsets. Conversely, Shan (\texttt{shn\_Mymr}) in Myanmar script is the lowest aligned language in both experiments, yielding a $\mu_{\text{Max}}$ score of only $0.06$ on FLORES and $0.0097$ on the Bible corpus.

Apertus 8B performs better on the Bible corpus than on FLORES,
indicating stronger low-resource language coverage.

\begin{table*}[t]
\centering
\footnotesize

\caption{
MEXA alignment scores
($\mu_{\text{Max}}$ / $\mu_{\text{Mean}}$)
comparing our model runs against reference baselines.
``---'' denotes configurations not yet run.
}

\label{tab:mexa_results}

\begin{adjustbox}{width=\textwidth}

\begin{tabular}{l|cc|cc|cc|cc|cc}

\toprule

\textbf{Model}
& \multicolumn{2}{c|}{\textbf{FLORES T1}}
& \multicolumn{2}{c|}{\textbf{FLORES T1 (2000)}}
& \multicolumn{2}{c|}{\textbf{Bible T1}}
& \multicolumn{2}{c|}{\textbf{FLORES Full}}
& \multicolumn{2}{c}{\textbf{Bible Full}} \\

& \multicolumn{2}{c|}{116 lang · 100 sent}
& \multicolumn{2}{c|}{116 lang · 2000 sent}
& \multicolumn{2}{c|}{101 lang · sPBC}
& \multicolumn{2}{c|}{204 lang · 2000 sent}
& \multicolumn{2}{c}{sPBC · all lang} \\

& $\mu_{\text{Max}}$
& $\mu_{\text{Mean}}$
& $\mu_{\text{Max}}$
& $\mu_{\text{Mean}}$
& $\mu_{\text{Max}}$
& $\mu_{\text{Mean}}$
& $\mu_{\text{Max}}$
& $\mu_{\text{Mean}}$
& $\mu_{\text{Max}}$
& $\mu_{\text{Mean}}$ \\

\midrule

\multicolumn{11}{l}{
\textit{Paper reference values \citep{kargaran2025mexa}:}
} \\

Llama 3.1 8B
& 0.6538 & 0.3963
& --- & ---
& 0.4212 & 0.2103
& --- & ---
& --- & --- \\

Mistral 7B v0.3
& 0.4716 & 0.2642
& --- & ---
& 0.2606 & 0.1198
& --- & ---
& --- & --- \\

\midrule

\multicolumn{11}{l}{
\textit{Our replications and extensions:}
} \\

Llama 3.1 8B
& 0.6735 & 0.4196
& 0.6065 & 0.3370
& 0.4180 & 0.2076
& \textbf{0.5611} & \textbf{0.3235}
& \textbf{0.0781} & \textbf{0.0320} \\

Mistral 7B v0.3
& 0.4980 & 0.2878
& 0.4066 & 0.2127
& 0.2571 & 0.1179
& 0.4102 & 0.2232
& 0.0465 & 0.0181 \\

Qwen3 8B Base
& 0.5759 & 0.3211
& 0.4838 & 0.2416
& 0.2970 & 0.1350
& 0.4723 & 0.2509
& 0.0499 & 0.0186 \\

Qwen3.5 9B Base
& \textbf{0.7814} & \textbf{0.5557}
& \textbf{0.7193} & \textbf{0.4748}
& \textbf{0.4821} & \textbf{0.2624}
& \textbf{0.5901} & \textbf{0.3727}
& \textbf{0.0845} & \textbf{0.0382} \\

Apertus 8B
& 0.3873 & 0.1637
& 0.2827 & 0.1176
& 0.4299 & 0.1896
& 0.3264 & 0.1267
& 0.0667 & 0.0237 \\

\bottomrule

\end{tabular}

\end{adjustbox}

\end{table*}


\begin{table*}[t]
\centering
\footnotesize

\caption{
MEXA alignment scores ($\mu_{\text{Max}}$ / $\mu_{\text{Mean}}$) comparing different model sizes in the Qwen3 family.
``---'' denotes configurations not run.
}

\label{tab:qwen_scaling}

\begin{adjustbox}{width=\textwidth}

\begin{tabular}{l|cc|cc|cc|cc|cc}

\toprule

\textbf{Model}
& \multicolumn{2}{c|}{\textbf{FLORES T1}}
& \multicolumn{2}{c|}{\textbf{FLORES T1 (2000)}}
& \multicolumn{2}{c|}{\textbf{Bible T1}}
& \multicolumn{2}{c|}{\textbf{FLORES Full}}
& \multicolumn{2}{c}{\textbf{Bible Full}} \\

& \multicolumn{2}{c|}{116 lang · 100 sent}
& \multicolumn{2}{c|}{116 lang · 2000 sent}
& \multicolumn{2}{c|}{101 lang · sPBC}
& \multicolumn{2}{c|}{204 lang · 2000 sent}
& \multicolumn{2}{c}{sPBC · all lang} \\

& $\mu_{\text{Max}}$
& $\mu_{\text{Mean}}$
& $\mu_{\text{Max}}$
& $\mu_{\text{Mean}}$
& $\mu_{\text{Max}}$
& $\mu_{\text{Mean}}$
& $\mu_{\text{Max}}$
& $\mu_{\text{Mean}}$
& $\mu_{\text{Max}}$
& $\mu_{\text{Mean}}$ \\

\midrule

Qwen3.5 9B Base
& \textbf{0.7814} & \textbf{0.5557}
& \textbf{0.7193} & \textbf{0.4748}
& \textbf{0.4821} & \textbf{0.2624}
& \textbf{0.5901} & \textbf{0.3727}
& \textbf{0.0845} & \textbf{0.0382} \\

Qwen3 8B Base
& 0.5759 & 0.3211
& 0.4838 & 0.2416
& 0.2970 & 0.1350
& \textbf{0.4723} & \textbf{0.2509}
& \textbf{0.0499} & \textbf{0.0186} \\

Qwen3 4B
& 0.4433 & 0.2327
& 0.3506 & 0.1658
& 0.1981 & 0.0918
& 0.2633 & 0.1200
& 0.0315 & 0.0120 \\

Qwen3 1.7B
& 0.4946 & 0.2158
& 0.4136 & 0.1531
& 0.1541 & 0.0612
& 0.3042 & 0.1106
& 0.0221 & 0.0072 \\

Qwen3 0.6B
& 0.3518 & 0.1504
& 0.2600 & 0.0936
& 0.1340 & 0.0466
& 0.1815 & 0.0639
& 0.0235 & 0.0067 \\

\bottomrule

\end{tabular}

\end{adjustbox}

\end{table*}

\subsection{Interactive Evaluation Dashboard}
\label{subsec:dashboard}

To facilitate a comprehensive exploration of the alignment results, I developed an interactive web visualization dashboard using React, Vite, and Plotly.js. The dashboard acts as an interface , providing a multi-dimensional overview of the experimental findings:

\begin{itemize}
    \item \textbf{Global Overview}: A comparative summary of the peak ($\mu_{\text{Max}}$) and average ($\mu_{\text{Mean}}$) alignment scores across all evaluated models and dataset variants (FLORES-200 and Bible).
    \item \textbf{Model Findings}: Dedicated pages for each model highlighting key statistics, such as the highest and lowest aligned languages, script-level performance differences, and complete score rankings.
    \item \textbf{Interactive Heatmaps}: A cell-grid visualization mapping MEXA alignment scores across 110+ languages and different model families.
    \item \textbf{2D/3D PCA Projections}: Dynamic dimensionality reduction plots visualizing the spatial distribution and clustering of sentence embeddings across languages, colored by writing system and MEXA score.
    \item \textbf{Cross-Model Correlation}: Interactive Pearson correlation matrices representing the similarity of alignment profiles across different model families.
\end{itemize}

This dashboard allows  to inspect fine-grained representation dynamics without having to parse raw data logs, to link the high-level alignment scores and qualitative language-script analyses.


% ============================================================
% Bibliography
% ============================================================

\bibliographystyle{plainnat}
\bibliography{references}

\end{document}